% ----- auto-generated from week_02/Week*.html via pandoc -----
\setchapterimage[6cm]{}
\setchapterpreamble[u]{\margintoc}
\chapter{Model Predictive Control --- Introduction}
\labch{week_02}

\section{Week 2: Model Predictive Control (MPC)}\label{week-2-model-predictive-control-mpc}}

From Receding-Horizon Optimisation to Real-Time Robotic Control

Control for Robots \textbar{} University of West London\\
AY 2025--26 \textbar{} 3-Hour Lecture + Activity Session

\subsection{Week 1 Recall: Computed Torque Control --- Exam Practice}\label{week-1-recall-computed-torque-control-exam-practice}}

Before we begin new material, you must demonstrate command of the Week 1 content.
Take out \textbf{pen and paper}. You have \textbf{20 minutes} to answer both questions below by hand, showing every step.
These are representative of the questions you will face in the final exam.

\paragraph{Question 1 --- Computed Torque Derivation (10 marks)}\label{question-1-computed-torque-derivation-10-marks}}

Given a 2-DOF planar manipulator with Euler--Lagrange dynamics

\textbackslash{[}
M(q)\textbackslash,\textbackslash ddot\{q\} + C(q,\textbackslash dot\{q\})\textbackslash,\textbackslash dot\{q\} + G(q) = \textbackslash tau
\textbackslash{]}

derive the computed torque control law

\textbackslash{[}
\textbackslash tau = M(q)\textbackslash bigl{[}\textbackslash ddot\{q\}\_d + K\_D\textbackslash,\textbackslash dot\{e\} + K\_P\textbackslash,e\textbackslash bigr{]} + C(q,\textbackslash dot\{q\})\textbackslash,\textbackslash dot\{q\} + G(q)
\textbackslash{]}

where \textbackslash(e = q\_d - q\textbackslash) is the tracking error.

\textbf{Step-by-step requirements:}

\begin{enumerate}
\tightlist
\item
  Substitute the control law \textbackslash(\textbackslash tau\textbackslash) into the equation of motion.
\item
  Cancel the Coriolis and gravity terms on both sides.
\item
  Left-multiply both sides by \textbackslash(M\^{}\{-1\}(q)\textbackslash) and justify why \textbackslash(M(q)\textbackslash) must be invertible (state the positive-definiteness property).
\item
  Rearrange using \textbackslash(\textbackslash ddot\{e\} = \textbackslash ddot\{q\}\_d - \textbackslash ddot\{q\}\textbackslash) to obtain the closed-loop error dynamics:
\end{enumerate}

\textbackslash{[}
\textbackslash ddot\{e\} + K\_D\textbackslash,\textbackslash dot\{e\} + K\_P\textbackslash,e = 0
\textbackslash{]}

Explain the physical significance of this result: what kind of system does it represent, and what does it guarantee about the tracking error?

\paragraph{Question 2 --- Lyapunov Stability Proof (10 marks)}\label{question-2-lyapunov-stability-proof-10-marks}}

For the computed torque controller derived in Question 1, construct the Lyapunov candidate function:

\textbackslash{[}
V(e, \textbackslash dot\{e\}) = \textbackslash frac\{1\}\{2\}\textbackslash,\textbackslash dot\{e\}\^{}T\textbackslash,M(q)\textbackslash,\textbackslash dot\{e\} \textbackslash;+\textbackslash; \textbackslash frac\{1\}\{2\}\textbackslash,e\^{}T\textbackslash,K\_P\textbackslash,e
\textbackslash{]}

\begin{enumerate}
\tightlist
\item
  State why \textbackslash(V \textgreater{} 0\textbackslash) for all \textbackslash((e, \textbackslash dot\{e\}) \textbackslash neq (0, 0)\textbackslash).
\item
  Differentiate \textbackslash(V\textbackslash) with respect to time. Show every term in \textbackslash(\textbackslash dot\{V\}\textbackslash).
\item
  Substitute the closed-loop dynamics \textbackslash(M(q)\textbackslash,\textbackslash ddot\{e\} = -C(q,\textbackslash dot\{q\})\textbackslash,\textbackslash dot\{e\} - K\_D\textbackslash,\textbackslash dot\{e\} - K\_P\textbackslash,e\textbackslash).
\item
  Apply the skew-symmetry property: \textbackslash(\textbackslash dot\{e\}\^{}T\textbackslash bigl{[}\textbackslash dot\{M\}(q) - 2C(q,\textbackslash dot\{q\})\textbackslash bigr{]}\textbackslash dot\{e\} = 0\textbackslash).
\item
  Show that \textbackslash(\textbackslash dot\{V\} = -\textbackslash dot\{e\}\^{}T\textbackslash,K\_D\textbackslash,\textbackslash dot\{e\} \textbackslash leq 0\textbackslash).
\item
  Invoke LaSalle\textquotesingle s invariance principle to conclude asymptotic stability of the equilibrium \textbackslash((e, \textbackslash dot\{e\}) = (0, 0)\textbackslash). State the invariance argument explicitly.
\end{enumerate}

\textbf{You have 20 minutes.} If you can reproduce both derivations cleanly, you will be able to do it in the exam. If you struggle, revise your Week 1 notes tonight. These two derivations form the foundation for everything that follows in this module.

\subsection{Learning Outcomes}\label{learning-outcomes}}

\textbf{By the end of this session, you will be able to:}

\begin{itemize}
\tightlist
\item
  \textbf{Formulate} MPC optimisation problems for robotic systems with state and control constraints
\item
  \textbf{Derive} the MPC control law from quadratic cost minimisation over a finite prediction horizon
\item
  \textbf{Implement} MPC for a 2-DOF manipulator, a differential-drive mobile robot, and a quadrotor drone
\item
  \textbf{Tune} the prediction horizon \textbackslash(N\textbackslash), state weight \textbackslash(Q\textbackslash), and control weight \textbackslash(R\textbackslash), and understand their effects on tracking performance, control effort, and computational cost
\item
  \textbf{Simulate} MPC controllers in MATLAB and interpret the resulting trajectory and error plots
\item
  \textbf{Compare} MPC with computed torque control in terms of model requirements, constraint handling, and computational demand
\end{itemize}

\subsection{Session Roadmap (3 Hours)}\label{session-roadmap-3-hours}}

0:00 -- 0:20

\paragraph{Previous Week Exam Practice}\label{previous-week-exam-practice}}

\begin{itemize}
\tightlist
\item
  Computed torque derivation --- pen and paper
\item
  Lyapunov stability proof --- pen and paper
\end{itemize}

0:20 -- 0:50

\paragraph{Part 1 --- MPC Theory \& Derivation}\label{part-1-mpc-theory-derivation}}

\begin{itemize}
\tightlist
\item
  What is MPC? The predict--optimise--apply paradigm
\item
  Optimisation problem formulation
\item
  State prediction over the horizon
\item
  Quadratic cost and the optimal control law
\item
  Receding horizon strategy
\end{itemize}

0:50 -- 1:20

\paragraph{Part 2 --- MPC for 2-DOF Manipulator}\label{part-2-mpc-for-2-dof-manipulator}}

\begin{itemize}
\tightlist
\item
  Linearisation and discretisation of manipulator dynamics
\item
  MPC formulation with 4 states and 2 controls
\item
  Full derivation and MATLAB code
\end{itemize}

1:20 -- 1:30

\paragraph{Break}\label{break}}

\begin{itemize}
\tightlist
\item
  10-minute break
\end{itemize}

1:30 -- 2:00

\paragraph{Part 3 --- MPC for Mobile Robot}\label{part-3-mpc-for-mobile-robot}}

\begin{itemize}
\tightlist
\item
  Differential-drive kinematic model
\item
  Linearisation around reference trajectory
\item
  MPC with velocity and turning-rate constraints
\item
  Full derivation and MATLAB code
\end{itemize}

2:00 -- 2:30

\paragraph{Part 4 --- MPC for Quadrotor Drone}\label{part-4-mpc-for-quadrotor-drone}}

\begin{itemize}
\tightlist
\item
  12-state quadrotor model linearised around hover
\item
  MPC with thrust and torque constraints
\item
  Full derivation and MATLAB code
\end{itemize}

2:30 -- 3:00

\paragraph{Part 5 --- MATLAB Implementation \& Activities}\label{part-5-matlab-implementation-activities}}

\begin{itemize}
\tightlist
\item
  Activity 1: Horizon length tuning experiment
\item
  Activity 2: Side-by-side comparison of all three systems
\item
  Think--Pair--Share and summary
\end{itemize}

\subsection{Part 1 --- MPC Theory \& Derivation}\label{part-1-mpc-theory-derivation-1}}

\subsubsection{1.1 What Is Model Predictive Control?}\label{what-is-model-predictive-control}}

Model Predictive Control (MPC) is an optimisation-based control strategy that operates on a simple but powerful principle:
at every time step, \textbf{predict} the future behaviour of the system over a finite horizon, \textbf{optimise} a control sequence to minimise a cost function,
and \textbf{apply} only the first control action before repeating the process at the next time step.

\paragraph{1. Predict}\label{predict}}

Use the system model \textbackslash(x\_\{k+1\} = A\textbackslash,x\_k + B\textbackslash,u\_k\textbackslash) to propagate the state forward \textbackslash(N\textbackslash) steps into the future, generating a predicted state trajectory \textbackslash(\textbackslash\{x\_\{k+1\}, x\_\{k+2\}, \textbackslash ldots, x\_\{k+N\}\textbackslash\}\textbackslash).

\paragraph{2. Optimise}\label{optimise}}

Find the control sequence \textbackslash(U = {[}u\_k, u\_\{k+1\}, \textbackslash ldots, u\_\{k+N-1\}{]}\^{}T\textbackslash) that minimises a cost function penalising state deviations from the reference and control effort, subject to constraints.

\paragraph{3. Apply}\label{apply}}

Apply only the first element \textbackslash(u\_k\^{}*\textbackslash) of the optimal sequence to the plant. At the next time step \textbackslash(k+1\textbackslash), measure the new state and repeat the entire process. This is the \textbf{receding horizon} strategy.

\emph{Figure 1: The MPC receding horizon concept. At each time step, the controller predicts future states over a finite horizon, optimises a control sequence, applies only the first input, then shifts the window forward.}

\subsubsection{1.2 The Optimisation Problem}\label{the-optimisation-problem}}

Consider a discrete-time linear time-invariant system:

\textbackslash{[} x\_\{k+1\} = A\textbackslash,x\_k + B\textbackslash,u\_k \textbackslash{]}

where \textbackslash(x\_k \textbackslash in \textbackslash mathbb\{R\}\^{}n\textbackslash) is the state, \textbackslash(u\_k \textbackslash in \textbackslash mathbb\{R\}\^{}m\textbackslash) is the control input, \textbackslash(A \textbackslash in \textbackslash mathbb\{R\}\^{}\{n \textbackslash times n\}\textbackslash) and \textbackslash(B \textbackslash in \textbackslash mathbb\{R\}\^{}\{n \textbackslash times m\}\textbackslash).

The MPC optimisation problem at time step \textbackslash(k\textbackslash) is:

\textbf{MPC Optimisation Problem}\\
\strut \\
\textbackslash{[}
\textbackslash min\_\{U\} \textbackslash; J = \textbackslash sum\_\{i=1\}\^{}\{N\} \textbackslash bigl(x\_\{k+i\} - x\_\{\textbackslash text\{ref\}\}\textbackslash bigr)\^{}T Q \textbackslash bigl(x\_\{k+i\} - x\_\{\textbackslash text\{ref\}\}\textbackslash bigr) \textbackslash;+\textbackslash; \textbackslash sum\_\{i=0\}\^{}\{N-1\} u\_\{k+i\}\^{}T\textbackslash,R\textbackslash,u\_\{k+i\}
\textbackslash{]}

subject to: \textbackslash(x\_\{k+i+1\} = A\textbackslash,x\_\{k+i\} + B\textbackslash,u\_\{k+i\}\textbackslash), \textbackslash(\textbackslash;u\_\{\textbackslash min\} \textbackslash leq u\_\{k+i\} \textbackslash leq u\_\{\textbackslash max\}\textbackslash), \textbackslash(\textbackslash;x\_\{\textbackslash min\} \textbackslash leq x\_\{k+i\} \textbackslash leq x\_\{\textbackslash max\}\textbackslash)

Here \textbackslash(Q \textbackslash succeq 0\textbackslash) is the state cost matrix (penalises deviation from the reference),
\textbackslash(R \textbackslash succ 0\textbackslash) is the control cost matrix (penalises large control inputs),
and \textbackslash(N\textbackslash) is the prediction horizon length.

\emph{Figure 2: MPC closed-loop control architecture. The optimizer uses a system model and constraints to compute the optimal control sequence; only the first input is applied to the plant. State feedback closes the loop.}

\subsubsection{1.3 State Prediction over the Horizon}\label{state-prediction-over-the-horizon}}

By recursively applying the dynamics, we can express all future states in terms of the current state \textbackslash(x\_k\textbackslash) and the control sequence \textbackslash(U\textbackslash):

\paragraph{Recursive State Prediction}\label{recursive-state-prediction}}

\textbf{Step 1:} One step ahead:

\textbackslash{[} x\_\{k+1\} = A\textbackslash,x\_k + B\textbackslash,u\_k \textbackslash{]}

\textbf{Step 2:} Two steps ahead:

\textbackslash{[} x\_\{k+2\} = A\textbackslash,x\_\{k+1\} + B\textbackslash,u\_\{k+1\} = A\^{}2\textbackslash,x\_k + A\textbackslash,B\textbackslash,u\_k + B\textbackslash,u\_\{k+1\} \textbackslash{]}

\textbf{Step 3:} General \textbackslash(i\textbackslash) steps ahead:

\textbackslash{[} x\_\{k+i\} = A\^{}i\textbackslash,x\_k + \textbackslash sum\_\{j=0\}\^{}\{i-1\} A\^{}\{i-1-j\}\textbackslash,B\textbackslash,u\_\{k+j\} \textbackslash{]}

\textbf{Step 4:} Stack all predictions into a single matrix equation. Define:

\textbackslash{[}
\textbackslash mathbf\{X\} = \textbackslash begin\{bmatrix\} x\_\{k+1\} \textbackslash\textbackslash{} x\_\{k+2\} \textbackslash\textbackslash{} \textbackslash vdots \textbackslash\textbackslash{} x\_\{k+N\} \textbackslash end\{bmatrix\}, \textbackslash quad
\textbackslash mathbf\{U\} = \textbackslash begin\{bmatrix\} u\_k \textbackslash\textbackslash{} u\_\{k+1\} \textbackslash\textbackslash{} \textbackslash vdots \textbackslash\textbackslash{} u\_\{k+N-1\} \textbackslash end\{bmatrix\}
\textbackslash{]}

Then:

\textbackslash{[}
\textbackslash mathbf\{X\} = \textbackslash mathcal\{A\}\textbackslash,x\_k + \textbackslash mathcal\{B\}\textbackslash,\textbackslash mathbf\{U\}
\textbackslash{]}

where:

\textbackslash{[}
\textbackslash mathcal\{A\} = \textbackslash begin\{bmatrix\} A \textbackslash\textbackslash{} A\^{}2 \textbackslash\textbackslash{} A\^{}3 \textbackslash\textbackslash{} \textbackslash vdots \textbackslash\textbackslash{} A\^{}N \textbackslash end\{bmatrix\}, \textbackslash quad
\textbackslash mathcal\{B\} = \textbackslash begin\{bmatrix\}
B \& 0 \& 0 \& \textbackslash cdots \& 0 \textbackslash\textbackslash{}
AB \& B \& 0 \& \textbackslash cdots \& 0 \textbackslash\textbackslash{}
A\^{}2B \& AB \& B \& \textbackslash cdots \& 0 \textbackslash\textbackslash{}
\textbackslash vdots \& \textbackslash vdots \& \textbackslash vdots \& \textbackslash ddots \& \textbackslash vdots \textbackslash\textbackslash{}
A\^{}\{N-1\}B \& A\^{}\{N-2\}B \& A\^{}\{N-3\}B \& \textbackslash cdots \& B
\textbackslash end\{bmatrix\}
\textbackslash{]}

\emph{Figure 3: State prediction visualization. The optimizer evaluates multiple control sequences (grey dashed) and selects the one minimizing the cost (green). Only the first control action u\_k* is applied; the rest are discarded and re-computed at the next step.}

\subsubsection{1.4 Quadratic Cost in Matrix Form}\label{quadratic-cost-in-matrix-form}}

Define the stacked reference \textbackslash(\textbackslash mathbf\{X\}\_\{\textbackslash text\{ref\}\} = {[}x\_\{\textbackslash text\{ref\}\}\^{}T, x\_\{\textbackslash text\{ref\}\}\^{}T, \textbackslash ldots, x\_\{\textbackslash text\{ref\}\}\^{}T{]}\^{}T \textbackslash in \textbackslash mathbb\{R\}\^{}\{Nn\}\textbackslash)
and the block-diagonal weight matrices:

\textbackslash{[}
\textbackslash mathbf\{Q\}\_N = \textbackslash text\{blkdiag\}(Q, Q, \textbackslash ldots, Q) \textbackslash in \textbackslash mathbb\{R\}\^{}\{Nn \textbackslash times Nn\}, \textbackslash quad
\textbackslash mathbf\{R\}\_N = \textbackslash text\{blkdiag\}(R, R, \textbackslash ldots, R) \textbackslash in \textbackslash mathbb\{R\}\^{}\{Nm \textbackslash times Nm\}
\textbackslash{]}

Substituting \textbackslash(\textbackslash mathbf\{X\} = \textbackslash mathcal\{A\}\textbackslash,x\_k + \textbackslash mathcal\{B\}\textbackslash,\textbackslash mathbf\{U\}\textbackslash) into the cost function:

\paragraph{Cost Function Expansion}\label{cost-function-expansion}}

\textbackslash{[}
J = (\textbackslash mathbf\{X\} - \textbackslash mathbf\{X\}\_\{\textbackslash text\{ref\}\})\^{}T \textbackslash mathbf\{Q\}\_N (\textbackslash mathbf\{X\} - \textbackslash mathbf\{X\}\_\{\textbackslash text\{ref\}\}) + \textbackslash mathbf\{U\}\^{}T \textbackslash mathbf\{R\}\_N \textbackslash mathbf\{U\}
\textbackslash{]}

Let \textbackslash(\textbackslash mathbf\{E\} = \textbackslash mathbf\{X\}\_\{\textbackslash text\{ref\}\} - \textbackslash mathcal\{A\}\textbackslash,x\_k\textbackslash) (the "free-response error"). Then \textbackslash(\textbackslash mathbf\{X\} - \textbackslash mathbf\{X\}\_\{\textbackslash text\{ref\}\} = \textbackslash mathcal\{B\}\textbackslash,\textbackslash mathbf\{U\} - \textbackslash mathbf\{E\}\textbackslash), so:

\textbackslash{[}
J = (\textbackslash mathcal\{B\}\textbackslash,\textbackslash mathbf\{U\} - \textbackslash mathbf\{E\})\^{}T \textbackslash mathbf\{Q\}\_N (\textbackslash mathcal\{B\}\textbackslash,\textbackslash mathbf\{U\} - \textbackslash mathbf\{E\}) + \textbackslash mathbf\{U\}\^{}T \textbackslash mathbf\{R\}\_N \textbackslash mathbf\{U\}
\textbackslash{]}

Expanding:

\textbackslash{[}
J = \textbackslash mathbf\{U\}\^{}T \textbackslash underbrace\{(\textbackslash mathcal\{B\}\^{}T \textbackslash mathbf\{Q\}\_N \textbackslash mathcal\{B\} + \textbackslash mathbf\{R\}\_N)\}\_\{H\} \textbackslash mathbf\{U\}
- 2\textbackslash,\textbackslash mathbf\{E\}\^{}T \textbackslash mathbf\{Q\}\_N \textbackslash mathcal\{B\}\textbackslash,\textbackslash mathbf\{U\}
+ \textbackslash mathbf\{E\}\^{}T \textbackslash mathbf\{Q\}\_N \textbackslash mathbf\{E\}
\textbackslash{]}

This is a strictly convex quadratic in \textbackslash(\textbackslash mathbf\{U\}\textbackslash) (since \textbackslash(H \textbackslash succ 0\textbackslash) when \textbackslash(R \textbackslash succ 0\textbackslash)), so the unique minimum is found by setting \textbackslash(\textbackslash frac\{\textbackslash partial J\}\{\textbackslash partial \textbackslash mathbf\{U\}\} = 0\textbackslash):

\textbackslash{[}
2\textbackslash,H\textbackslash,\textbackslash mathbf\{U\}\^{}* - 2\textbackslash,\textbackslash mathcal\{B\}\^{}T \textbackslash mathbf\{Q\}\_N \textbackslash mathbf\{E\} = 0
\textbackslash{]}

\subsubsection{1.5 The Optimal Control Law}\label{the-optimal-control-law}}

\textbf{MPC Optimal Control Sequence}\\
\strut \\
\textbackslash{[}
\textbackslash mathbf\{U\}\^{}* = \textbackslash bigl(\textbackslash mathcal\{B\}\^{}T \textbackslash mathbf\{Q\}\_N \textbackslash mathcal\{B\} + \textbackslash mathbf\{R\}\_N\textbackslash bigr)\^{}\{-1\} \textbackslash mathcal\{B\}\^{}T \textbackslash mathbf\{Q\}\_N \textbackslash bigl(\textbackslash mathbf\{X\}\_\{\textbackslash text\{ref\}\} - \textbackslash mathcal\{A\}\textbackslash,x\_k\textbackslash bigr)
\textbackslash{]}

The matrix \textbackslash(H = \textbackslash mathcal\{B\}\^{}T \textbackslash mathbf\{Q\}\_N \textbackslash mathcal\{B\} + \textbackslash mathbf\{R\}\_N\textbackslash) is of dimension \textbackslash(Nm \textbackslash times Nm\textbackslash). It must be inverted (or factorised) at every time step. This is the primary computational cost of MPC.

\subsubsection{1.6 Receding Horizon: Apply Only the First Input}\label{receding-horizon-apply-only-the-first-input}}

The full optimal sequence is \textbackslash(\textbackslash mathbf\{U\}\^{}* = {[}u\_k\^{}\{*T\}, u\_\{k+1\}\^{}\{*T\}, \textbackslash ldots, u\_\{k+N-1\}\^{}\{*T\}{]}\^{}T\textbackslash).
In the receding-horizon strategy, we \textbf{apply only the first element}:

\textbackslash{[} u\_k = u\_k\^{}* = \textbackslash text\{first \} m \textbackslash text\{ entries of \} \textbackslash mathbf\{U\}\^{}* \textbackslash{]}

At the next time step, the state \textbackslash(x\_\{k+1\}\textbackslash) is measured (or estimated), and the entire optimisation is re-solved.
This feedback mechanism gives MPC its robustness to disturbances and model errors.

\emph{Figure 4: MPC flowchart showing the predict-optimise-apply loop. The key insight is that only the first control input is applied (highlighted in green), and the entire optimisation is re-solved at every time step.}

\paragraph{Think--Pair--Share (3 min)}\label{thinkpairshare-3-min}}

\textbf{Question:} Why does MPC re-solve the optimisation at every time step instead of applying the entire precomputed sequence \textbackslash(\textbackslash mathbf\{U\}\^{}*\textbackslash)?

Hint: Think about what happens if the model is imperfect, or if a disturbance pushes the state away from the predicted trajectory.
Compare with open-loop vs closed-loop control.

PSEUDOCODE --- MPC Receding Horizon Algorithm

INITIALISE: state x\_0, horizon N, weights Q and R
BUILD: prediction matrices A\_cal and B\_cal from (A, B, N)
BUILD: block-diagonal weights Q\_N = blkdiag(Q,...,Q), R\_N = blkdiag(R,...,R)
COMPUTE: H = B\_cal\textquotesingle{} * Q\_N * B\_cal + R\_N
COMPUTE: H\_inv = inv(H) (do this ONCE if A, B are constant)
FOR k = 0, 1, 2, ... DO
MEASURE (or estimate) current state x\_k
BUILD reference vector: X\_ref = {[}x\_ref; x\_ref; ...; x\_ref{]} (N copies)
COMPUTE free-response error: E = X\_ref - A\_cal * x\_k
COMPUTE optimal sequence: U\_star = H\_inv * B\_cal\textquotesingle{} * Q\_N * E
EXTRACT first control: u\_k = U\_star(1:m)
APPLY u\_k to the plant
ADVANCE to next time step
END FOR

\subsubsection{1.7 Effect of Tuning Parameters}\label{effect-of-tuning-parameters}}

\begin{longtable}[]{@{}lll@{}}
\toprule\noalign{}
Parameter & Increase Effect & Decrease Effect \\
\midrule\noalign{}
\endhead
\bottomrule\noalign{}
\endlastfoot
Prediction horizon \textbackslash(N\textbackslash) & Better anticipation of future references; more stable; higher computational cost (\textbackslash(O(N\^{}3 m\^{}3)\textbackslash)) & More myopic; may oscillate or become unstable; cheaper to compute \\
State weight \textbackslash(Q\textbackslash) & Tighter tracking; more aggressive control; risk of actuator saturation & Relaxed tracking; smoother control inputs; larger steady-state error \\
Control weight \textbackslash(R\textbackslash) & Penalises large inputs; smoother control; slower response; energy-efficient & Allows large inputs; fast response; risk of actuator damage \\
\end{longtable}

\textbf{Practical Guideline:} Start with \textbackslash(Q = I\textbackslash), \textbackslash(R = 0.1 I\textbackslash), and \textbackslash(N = 10\textbackslash). Then adjust: increase \textbackslash(Q\textbackslash) if tracking is too loose, increase \textbackslash(R\textbackslash) if control inputs saturate, and increase \textbackslash(N\textbackslash) until performance no longer improves significantly.

\subsection{Part 2 --- MPC for 2-DOF Manipulator}\label{part-2-mpc-for-2-dof-manipulator-1}}

\subsubsection{2.1 The Nonlinear Dynamics}\label{the-nonlinear-dynamics}}

Recall from Week 1 the 2-DOF planar manipulator dynamics:

\textbackslash{[}
M(q)\textbackslash,\textbackslash ddot\{q\} + C(q,\textbackslash dot\{q\})\textbackslash,\textbackslash dot\{q\} + G(q) = \textbackslash tau
\textbackslash{]}

with state vector \textbackslash(x = {[}q\_1, q\_2, \textbackslash dot\{q\}\_1, \textbackslash dot\{q\}\_2{]}\^{}T \textbackslash in \textbackslash mathbb\{R\}\^{}4\textbackslash) and control input \textbackslash(u = {[}\textbackslash tau\_1, \textbackslash tau\_2{]}\^{}T \textbackslash in \textbackslash mathbb\{R\}\^{}2\textbackslash).

\emph{Figure 5: 2-DOF planar manipulator with joint angles θ1, θ2, link lengths l1, l2, masses m1, m2, and torque inputs τ1, τ2. The coordinate frame is shown at the base with gravity acting downward.}

\subsubsection{2.2 Linearisation around an Operating Point}\label{linearisation-around-an-operating-point}}

MPC (in its basic linear form) requires a linear discrete-time model. We linearise the manipulator dynamics around a nominal operating point
\textbackslash((q\_0, \textbackslash dot\{q\}\_0) = (0, 0)\textbackslash) (the hanging-down equilibrium).

\paragraph{Linearisation Procedure}\label{linearisation-procedure}}

The continuous-time state-space form is:

\textbackslash{[}
\textbackslash dot\{x\} = f(x, u) = \textbackslash begin\{bmatrix\} \textbackslash dot\{q\} \textbackslash\textbackslash{} M\^{}\{-1\}(q)\textbackslash bigl{[}\textbackslash tau - C(q,\textbackslash dot\{q\})\textbackslash dot\{q\} - G(q)\textbackslash bigr{]} \textbackslash end\{bmatrix\}
\textbackslash{]}

At the operating point \textbackslash(x\_0 = {[}0, 0, 0, 0{]}\^{}T\textbackslash), \textbackslash(u\_0 = G(q\_0)\textbackslash):

\textbackslash{[}
A\_c = \textbackslash left.\textbackslash frac\{\textbackslash partial f\}\{\textbackslash partial x\}\textbackslash right\textbar\_\{x\_0, u\_0\}, \textbackslash quad
B\_c = \textbackslash left.\textbackslash frac\{\textbackslash partial f\}\{\textbackslash partial u\}\textbackslash right\textbar\_\{x\_0, u\_0\}
\textbackslash{]}

For the 2-DOF arm with \textbackslash(m\_1 = m\_2 = 1\textbackslash,\textbackslash text\{kg\}\textbackslash), \textbackslash(l\_1 = l\_2 = 1\textbackslash,\textbackslash text\{m\}\textbackslash), evaluated at \textbackslash(q = 0\textbackslash):

\textbackslash{[}
M(0) = \textbackslash begin\{bmatrix\} m\_1 l\_1\^{}2 + m\_2(l\_1\^{}2 + l\_2\^{}2 + 2l\_1 l\_2) \& m\_2(l\_2\^{}2 + l\_1 l\_2) \textbackslash\textbackslash{} m\_2(l\_2\^{}2 + l\_1 l\_2) \& m\_2 l\_2\^{}2 \textbackslash end\{bmatrix\}
= \textbackslash begin\{bmatrix\} 4 \& 2 \textbackslash\textbackslash{} 2 \& 1 \textbackslash end\{bmatrix\}
\textbackslash{]}

We evaluate at a slightly offset configuration \textbackslash(q\_0 = {[}0.1, 0.1{]}\^{}T\textbackslash) to ensure non-singularity:

\textbackslash{[}
M(q\_0) \textbackslash approx \textbackslash begin\{bmatrix\} 3.96 \& 1.99 \textbackslash\textbackslash{} 1.99 \& 1.0 \textbackslash end\{bmatrix\}, \textbackslash quad
M\^{}\{-1\}(q\_0) \textbackslash approx \textbackslash begin\{bmatrix\} 50.25 \& -100.0 \textbackslash\textbackslash{} -100.0 \& 199.0 \textbackslash end\{bmatrix\}
\textbackslash{]}

\subsubsection{2.3 Discretisation}\label{discretisation}}

\textbackslash{[}
A\_d = e\^{}\{A\_c \textbackslash Delta t\} \textbackslash approx I + A\_c \textbackslash Delta t, \textbackslash quad
B\_d \textbackslash approx B\_c \textbackslash Delta t
\textbackslash{]}

\subsubsection{2.4 MPC Design for the Manipulator}\label{mpc-design-for-the-manipulator}}

\paragraph{Weight Matrix Selection}\label{weight-matrix-selection}}

\textbackslash{[}
Q = \textbackslash text\{diag\}(100, 100, 1, 1), \textbackslash quad R = \textbackslash text\{diag\}(0.01, 0.01)
\textbackslash{]}

\subsubsection{2.5 MATLAB Implementation --- Manipulator MPC}\label{matlab-implementation-manipulator-mpc}}

\%\% === MPC for 2-DOF Manipulator ===
\% System parameters
params.m1 = 1; params.m2 = 1;
params.l1 = 1; params.l2 = 1;
params.g = 9.81;
\% Linearisation point
q0 = {[}0.1; 0.1{]}; dq0 = {[}0; 0{]};
\% Build M(q0), compute M\_inv
c2 = cos(q0(2));
M11 = params.m1*params.l1\^{}2 + params.m2*(params.l1\^{}2 + params.l2\^{}2 + 2*params.l1*params.l2*c2);
M12 = params.m2*(params.l2\^{}2 + params.l1*params.l2*c2);
M22 = params.m2*params.l2\^{}2;
M0 = {[}M11, M12; M12, M22{]};
M0\_inv = inv(M0);
\% Continuous-time A, B (linearised around q0, dq0=0)
Ac = {[}zeros(2), eye(2); zeros(2), zeros(2){]};
Bc = {[}zeros(2); M0\_inv{]};
\% Discretise (Euler, dt = 0.01 s)
dt = 0.01;
Ad = eye(4) + Ac*dt;
Bd = Bc*dt;
\% MPC parameters
N = 20; \% Prediction horizon
Q = diag({[}100, 100, 1, 1{]}); \% State weight
R = diag({[}0.01, 0.01{]}); \% Control weight
\% Build prediction matrices
n = 4; m = 2;
A\_cal = zeros(n*N, n);
B\_cal = zeros(n*N, m*N);
for i = 1:N
A\_cal((i-1)*n+1:i*n, :) = Ad\^{}i;
for j = 1:i
B\_cal((i-1)*n+1:i*n, (j-1)*m+1:j*m) = Ad\^{}(i-j) * Bd;
end
end
\% Block-diagonal weights
Q\_N = kron(eye(N), Q);
R\_N = kron(eye(N), R);
\% Hessian and its inverse (constant for LTI system)
H = B\_cal\textquotesingle{} * Q\_N * B\_cal + R\_N;
H\_inv = inv(H);
F = B\_cal\textquotesingle{} * Q\_N;
\% Simulation
t\_span = 0:dt:20;
x = {[}0; 0; 0; 0{]}; \% Initial state
x\_hist = zeros(4, length(t\_span));
u\_hist = zeros(2, length(t\_span));
x\_hist(:,1) = x;
for k = 1:length(t\_span)-1
t = t\_span(k);
\% Reference trajectory (sinusoidal)
x\_ref\_k = zeros(n*N, 1);
for i = 1:N
t\_future = t + i*dt;
x\_ref\_k((i-1)*n+1:i*n) = {[}0.5*sin(t\_future); 0.5*cos(t\_future); ...
0.5*cos(t\_future); -0.5*sin(t\_future){]};
end
\% MPC optimal control
E = x\_ref\_k - A\_cal * x;
U\_star = H\_inv * F * E;
u = U\_star(1:m); \% Apply first control only
\% Clamp torques (constraint enforcement)
u = max(min(u, 50), -50);
\% Simulate true nonlinear dynamics (one Euler step)
q = x(1:2); dq = x(3:4);
{[}M\_true, C\_true, G\_true{]} = manipulator\_model(q, dq, params);
ddq = M\_true \textbackslash{} (u - C\_true*dq - G\_true);
x = x + {[}dq; ddq{]}*dt;
x\_hist(:,k+1) = x;
u\_hist(:,k) = u;
end

\textbf{Note:} The controller uses a \emph{linear} model, but the plant simulation uses the \emph{full nonlinear} dynamics. This mismatch is realistic and tests the robustness of MPC.

\subsection{Part 3 --- MPC for Differential-Drive Mobile Robot}\label{part-3-mpc-for-differential-drive-mobile-robot}}

\subsubsection{3.1 Kinematic Model}\label{kinematic-model}}

A differential-drive mobile robot moves in the plane with state \textbackslash(x = {[}p\_x, p\_y, \textbackslash theta{]}\^{}T\textbackslash) and control inputs \textbackslash(u = {[}v, \textbackslash omega{]}\^{}T\textbackslash):

\textbf{Nonlinear Kinematic Model}\\
\strut \\
\textbackslash{[}
\textbackslash begin\{bmatrix\} \textbackslash dot\{p\}\_x \textbackslash\textbackslash{} \textbackslash dot\{p\}\_y \textbackslash\textbackslash{} \textbackslash dot\{\textbackslash theta\} \textbackslash end\{bmatrix\}
= \textbackslash begin\{bmatrix\} v\textbackslash,\textbackslash cos\textbackslash theta \textbackslash\textbackslash{} v\textbackslash,\textbackslash sin\textbackslash theta \textbackslash\textbackslash{} \textbackslash omega \textbackslash end\{bmatrix\}
\textbackslash{]}

\emph{Figure 6: Top-down view of a differential-drive mobile robot. The robot state is (p\_x, p\_y, θ), controlled by linear velocity v and angular velocity ω. The reference path is shown as a dashed curve.}

\subsubsection{3.2 Linearisation around a Reference Trajectory}\label{linearisation-around-a-reference-trajectory}}

\paragraph{Jacobian Linearisation}\label{jacobian-linearisation}}

\textbackslash{[}
A\_c = \textbackslash begin\{bmatrix\} 0 \& 0 \& -v\^{}r \textbackslash sin\textbackslash theta\^{}r \textbackslash\textbackslash{} 0 \& 0 \& v\^{}r \textbackslash cos\textbackslash theta\^{}r \textbackslash\textbackslash{} 0 \& 0 \& 0 \textbackslash end\{bmatrix\}, \textbackslash quad
B\_c = \textbackslash begin\{bmatrix\} \textbackslash cos\textbackslash theta\^{}r \& 0 \textbackslash\textbackslash{} \textbackslash sin\textbackslash theta\^{}r \& 0 \textbackslash\textbackslash{} 0 \& 1 \textbackslash end\{bmatrix\}
\textbackslash{]}

For \textbackslash(\textbackslash theta\^{}r = 0\textbackslash), \textbackslash(v\^{}r = 2.5\textbackslash,\textbackslash text\{m/s\}\textbackslash):

\textbackslash{[}
A\_c = \textbackslash begin\{bmatrix\} 0 \& 0 \& 0 \textbackslash\textbackslash{} 0 \& 0 \& 2.5 \textbackslash\textbackslash{} 0 \& 0 \& 0 \textbackslash end\{bmatrix\}, \textbackslash quad
B\_c = \textbackslash begin\{bmatrix\} 1 \& 0 \textbackslash\textbackslash{} 0 \& 0 \textbackslash\textbackslash{} 0 \& 1 \textbackslash end\{bmatrix\}
\textbackslash{]}

\subsubsection{3.3 Constraints}\label{constraints}}

\begin{longtable}[]{@{}lll@{}}
\toprule\noalign{}
Constraint & Physical Meaning & Value \\
\midrule\noalign{}
\endhead
\bottomrule\noalign{}
\endlastfoot
\textbackslash(0 \textbackslash leq v \textbackslash leq 5\textbackslash,\textbackslash text\{m/s\}\textbackslash) & No backward motion; maximum speed & {[}0, 5{]} \\
\textbackslash(-2 \textbackslash leq \textbackslash omega \textbackslash leq 2\textbackslash,\textbackslash text\{rad/s\}\textbackslash) & Maximum turning rate & {[}-2, 2{]} \\
\end{longtable}

\subsubsection{3.4 MPC Design}\label{mpc-design}}

\textbackslash{[} Q = \textbackslash text\{diag\}(50, 50, 10), \textbackslash quad R = \textbackslash text\{diag\}(0.1, 0.1) \textbackslash{]}

\subsubsection{3.5 MATLAB Implementation --- Mobile Robot MPC}\label{matlab-implementation-mobile-robot-mpc}}

\%\% === MPC for Differential-Drive Mobile Robot ===
v\_ref = 2.5; dt = 0.05;
Ac\_mob = {[}0, 0, 0; 0, 0, v\_ref; 0, 0, 0{]};
Bc\_mob = {[}1, 0; 0, 0; 0, 1{]};
Ad\_mob = eye(3) + Ac\_mob*dt;
Bd\_mob = Bc\_mob*dt;
N\_mob = 15; Q\_mob = diag({[}50, 50, 10{]}); R\_mob = diag({[}0.1, 0.1{]});
n\_mob = 3; m\_mob = 2;
\% Build prediction matrices (same structure as manipulator)
A\_cal\_mob = zeros(n\_mob*N\_mob, n\_mob);
B\_cal\_mob = zeros(n\_mob*N\_mob, m\_mob*N\_mob);
for i = 1:N\_mob
A\_cal\_mob((i-1)*n\_mob+1:i*n\_mob, :) = Ad\_mob\^{}i;
for j = 1:i
B\_cal\_mob((i-1)*n\_mob+1:i*n\_mob, (j-1)*m\_mob+1:j*m\_mob) = Ad\_mob\^{}(i-j)*Bd\_mob;
end
end
Q\_N\_mob = kron(eye(N\_mob), Q\_mob);
R\_N\_mob = kron(eye(N\_mob), R\_mob);
H\_mob = B\_cal\_mob\textquotesingle{} * Q\_N\_mob * B\_cal\_mob + R\_N\_mob;
H\_inv\_mob = inv(H\_mob);
F\_mob = B\_cal\_mob\textquotesingle{} * Q\_N\_mob;
\% Simulation (circular reference)
t\_span\_mob = 0:dt:20;
x\_mob = {[}0; 0; 0{]};
R\_circle = 5; omega\_circle = 0.5;
for k = 1:length(t\_span\_mob)-1
t = t\_span\_mob(k);
x\_ref\_mob = zeros(n\_mob*N\_mob, 1);
for i = 1:N\_mob
t\_f = t + i*dt;
x\_ref\_mob((i-1)*n\_mob+1:i*n\_mob) = {[}R\_circle*cos(omega\_circle*t\_f); ...
R\_circle*sin(omega\_circle*t\_f); ...
omega\_circle*t\_f + pi/2{]};
end
E\_mob = x\_ref\_mob - A\_cal\_mob * x\_mob;
U\_star\_mob = H\_inv\_mob * F\_mob * E\_mob;
u\_mob = U\_star\_mob(1:m\_mob);
u\_mob(1) = max(min(u\_mob(1), 5), 0);
u\_mob(2) = max(min(u\_mob(2), 2), -2);
theta = x\_mob(3);
x\_mob = x\_mob + {[}u\_mob(1)*cos(theta); u\_mob(1)*sin(theta); u\_mob(2){]}*dt;
end

\textbf{Linearisation Caveat:} The linearisation is performed at a single operating point. In practice, one re-linearises at each time step or uses nonlinear MPC.

\subsection{Part 4 --- MPC for Quadrotor Drone}\label{part-4-mpc-for-quadrotor-drone-1}}

\subsubsection{4.1 The 12-State Quadrotor Model}\label{the-12-state-quadrotor-model}}

A quadrotor has 6 degrees of freedom (3 translational + 3 rotational), leading to a 12-dimensional state vector:

\textbackslash{[}
x = \textbackslash begin\{bmatrix\} p\_x \textbackslash\textbackslash{} p\_y \textbackslash\textbackslash{} p\_z \textbackslash\textbackslash{} v\_x \textbackslash\textbackslash{} v\_y \textbackslash\textbackslash{} v\_z \textbackslash\textbackslash{} \textbackslash phi \textbackslash\textbackslash{} \textbackslash theta \textbackslash\textbackslash{} \textbackslash psi \textbackslash\textbackslash{} \textbackslash omega\_x \textbackslash\textbackslash{} \textbackslash omega\_y \textbackslash\textbackslash{} \textbackslash omega\_z \textbackslash end\{bmatrix\}
\textbackslash quad \textbackslash text\{(position, velocity, Euler angles, angular rates)\}
\textbackslash{]}

\emph{Figure 7: Quadrotor UAV schematic showing four rotors with alternating CW/CCW rotation directions, thrust forces F1-F4, body frame axes (x\_B, y\_B, z\_B), and Euler angles (roll φ, pitch θ, yaw ψ). The total thrust is T = F1+F2+F3+F4.}

The control inputs are thrust and three body-frame torques:

\textbackslash{[}
u = \textbackslash begin\{bmatrix\} T \textbackslash\textbackslash{} \textbackslash tau\_x \textbackslash\textbackslash{} \textbackslash tau\_y \textbackslash\textbackslash{} \textbackslash tau\_z \textbackslash end\{bmatrix\}
\textbackslash{]}

\subsubsection{4.2 Linearised Hover Dynamics}\label{linearised-hover-dynamics}}

\paragraph{Linearised Hover Dynamics}\label{linearised-hover-dynamics-1}}

At hover (\textbackslash(\textbackslash phi = \textbackslash theta = \textbackslash psi = 0\textbackslash), \textbackslash(T\_0 = mg\textbackslash)):

\textbackslash{[}
\textbackslash dot\{v\}\_x \textbackslash approx g\textbackslash,\textbackslash theta, \textbackslash quad \textbackslash dot\{v\}\_y \textbackslash approx -g\textbackslash,\textbackslash phi, \textbackslash quad \textbackslash dot\{v\}\_z \textbackslash approx \textbackslash frac\{T - mg\}\{m\}
\textbackslash{]}
\textbackslash{[}
\textbackslash dot\{\textbackslash omega\}\_x \textbackslash approx \textbackslash frac\{\textbackslash tau\_x\}\{I\_\{xx\}\}, \textbackslash quad \textbackslash dot\{\textbackslash omega\}\_y \textbackslash approx \textbackslash frac\{\textbackslash tau\_y\}\{I\_\{yy\}\}, \textbackslash quad \textbackslash dot\{\textbackslash omega\}\_z \textbackslash approx \textbackslash frac\{\textbackslash tau\_z\}\{I\_\{zz\}\}
\textbackslash{]}

\subsubsection{4.3 MPC Design for the Quadrotor}\label{mpc-design-for-the-quadrotor}}

\textbackslash{[}
Q = \textbackslash text\{diag\}(\textbackslash underbrace\{80, 80, 120\}\_\{\textbackslash text\{position\}\}, \textbackslash underbrace\{1, 1, 1\}\_\{\textbackslash text\{velocity\}\}, \textbackslash underbrace\{50, 50, 50\}\_\{\textbackslash text\{orientation\}\}, \textbackslash underbrace\{1, 1, 1\}\_\{\textbackslash text\{angular rate\}\})
\textbackslash{]}
\textbackslash{[}
R = \textbackslash text\{diag\}(0.1, 0.1, 0.1, 0.1)
\textbackslash{]}

\subsubsection{4.4 Constraints}\label{constraints-1}}

\begin{longtable}[]{@{}lll@{}}
\toprule\noalign{}
Constraint & Physical Meaning & Value \\
\midrule\noalign{}
\endhead
\bottomrule\noalign{}
\endlastfoot
\textbackslash(0 \textbackslash leq T \textbackslash leq 2mg\textbackslash) & Thrust range & {[}0, 19.62{]} N \\
\textbackslash(-2 \textbackslash leq \textbackslash tau\_x, \textbackslash tau\_y, \textbackslash tau\_z \textbackslash leq 2\textbackslash) & Maximum body torques & {[}-2, 2{]} N·m \\
\textbackslash(-30\^{}\textbackslash circ \textbackslash leq \textbackslash phi, \textbackslash theta \textbackslash leq 30\^{}\textbackslash circ\textbackslash) & Tilt angle limits & {[}-0.524, 0.524{]} rad \\
\end{longtable}

\subsubsection{4.5 MATLAB Implementation --- Quadrotor MPC}\label{matlab-implementation-quadrotor-mpc}}

\%\% === MPC for Quadrotor Drone ===
quad.m = 1.0; quad.Ixx = 0.1; quad.Iyy = 0.1; quad.Izz = 0.2; quad.g = 9.81;
dt\_q = 0.02; n\_q = 12; m\_q = 4;
Ac\_q = zeros(12,12);
Ac\_q(1:3, 4:6) = eye(3);
Ac\_q(4,8) = quad.g;
Ac\_q(5,7) = -quad.g;
Ac\_q(7:9, 10:12) = eye(3);
Bc\_q = zeros(12,4);
Bc\_q(6,1) = 1/quad.m;
Bc\_q(10,2) = 1/quad.Ixx;
Bc\_q(11,3) = 1/quad.Iyy;
Bc\_q(12,4) = 1/quad.Izz;
Ad\_q = eye(12) + Ac\_q*dt\_q;
Bd\_q = Bc\_q*dt\_q;
N\_q = 15;
Q\_q = diag({[}80, 80, 120, 1, 1, 1, 50, 50, 50, 1, 1, 1{]});
R\_q = diag({[}0.1, 0.1, 0.1, 0.1{]});
\% Build prediction matrices and solve (same pattern as above)
A\_cal\_q = zeros(n\_q*N\_q, n\_q);
B\_cal\_q = zeros(n\_q*N\_q, m\_q*N\_q);
for i = 1:N\_q
A\_cal\_q((i-1)*n\_q+1:i*n\_q, :) = Ad\_q\^{}i;
for j = 1:i
B\_cal\_q((i-1)*n\_q+1:i*n\_q, (j-1)*m\_q+1:j*m\_q) = Ad\_q\^{}(i-j)*Bd\_q;
end
end
H\_q = B\_cal\_q\textquotesingle{} * kron(eye(N\_q),Q\_q) * B\_cal\_q + kron(eye(N\_q),R\_q);
H\_inv\_q = inv(H\_q);
F\_q = B\_cal\_q\textquotesingle{} * kron(eye(N\_q),Q\_q);
\% Simulation loop follows same pattern with helical reference trajectory

\textbf{Key Observation:} The MPC controller for the quadrotor uses a 12-state linear model but drives the full nonlinear plant. The receding-horizon feedback mechanism compensates for linearisation errors at each step.

\subsection{Part 5 --- MATLAB Implementation \& Activities}\label{part-5-matlab-implementation-activities-1}}

\subsubsection{Activity 1: Prediction Horizon Tuning (10 min)}\label{activity-1-prediction-horizon-tuning-10-min}}

\paragraph{Task: Investigate the Effect of Horizon Length \textbackslash(N\textbackslash)}\label{task-investigate-the-effect-of-horizon-length-n}}

Run the manipulator MPC code with three different prediction horizons.

\begin{longtable}[]{@{}llll@{}}
\toprule\noalign{}
Horizon \textbackslash(N\textbackslash) & Peak Position Error & Steady-State Error & Computation Time \\
\midrule\noalign{}
\endhead
\bottomrule\noalign{}
\endlastfoot
\textbackslash(N = 5\textbackslash) & {[}record{]} & {[}record{]} & {[}record{]} \\
\textbackslash(N = 20\textbackslash) & {[}record{]} & {[}record{]} & {[}record{]} \\
\textbackslash(N = 50\textbackslash) & {[}record{]} & {[}record{]} & {[}record{]} \\
\end{longtable}

\subsubsection{Activity 2: Three-System Comparison (15 min)}\label{activity-2-three-system-comparison-15-min}}

\paragraph{Task: Run MPC for All Three Systems Side by Side}\label{task-run-mpc-for-all-three-systems-side-by-side}}

\begin{enumerate}
\tightlist
\item
  Compare the MPC tracking error with the PD tracking error for each system.
\item
  Examine the control input plots. Does MPC produce smoother control signals?
\item
  Look at the constraint enforcement: does PD ever violate the limits?
\end{enumerate}

\subsubsection{Activity 3: MPC vs Computed Torque Discussion (5 min)}\label{activity-3-mpc-vs-computed-torque-discussion-5-min}}

\paragraph{Think--Pair--Share}\label{thinkpairshare}}

\textbf{Question 1:} When would you prefer MPC over computed torque control?

\textbf{Question 2:} When would computed torque control be the better choice?

\subsection{Summary \& Key Takeaways}\label{summary-key-takeaways}}

\subsubsection{Core MPC Equations}\label{core-mpc-equations}}

\begin{longtable}[]{@{}ll@{}}
\toprule\noalign{}
Equation & Description \\
\midrule\noalign{}
\endhead
\bottomrule\noalign{}
\endlastfoot
\textbackslash(x\_\{k+1\} = A\textbackslash,x\_k + B\textbackslash,u\_k\textbackslash) & Discrete-time linear system model \\
\textbackslash(\textbackslash mathbf\{X\} = \textbackslash mathcal\{A\}\textbackslash,x\_k + \textbackslash mathcal\{B\}\textbackslash,\textbackslash mathbf\{U\}\textbackslash) & Stacked state prediction over horizon \textbackslash(N\textbackslash) \\
\textbackslash(J = (\textbackslash mathbf\{X\} - \textbackslash mathbf\{X\}\_\{\textbackslash text\{ref\}\})\^{}T \textbackslash mathbf\{Q\}\_N (\textbackslash mathbf\{X\} - \textbackslash mathbf\{X\}\_\{\textbackslash text\{ref\}\}) + \textbackslash mathbf\{U\}\^{}T \textbackslash mathbf\{R\}\_N \textbackslash mathbf\{U\}\textbackslash) & Quadratic cost function \\
\textbackslash(\textbackslash mathbf\{U\}\^{}* = (\textbackslash mathcal\{B\}\^{}T \textbackslash mathbf\{Q\}\_N \textbackslash mathcal\{B\} + \textbackslash mathbf\{R\}\_N)\^{}\{-1\} \textbackslash mathcal\{B\}\^{}T \textbackslash mathbf\{Q\}\_N (\textbackslash mathbf\{X\}\_\{\textbackslash text\{ref\}\} - \textbackslash mathcal\{A\}\textbackslash,x\_k)\textbackslash) & Optimal control law (unconstrained) \\
\textbackslash(u\_k = \textbackslash mathbf\{U\}\^{}*(1\{:\}m)\textbackslash) & Receding horizon: apply only first input \\
\end{longtable}

\subsubsection{MPC vs Computed Torque Control}\label{mpc-vs-computed-torque-control}}

\begin{longtable}[]{@{}lll@{}}
\toprule\noalign{}
Property & MPC & Computed Torque \\
\midrule\noalign{}
\endhead
\bottomrule\noalign{}
\endlastfoot
\textbf{Model Knowledge} & Requires linear (or nonlinear) model; can use approximate models & Requires full nonlinear dynamics \textbackslash(M(q)\textbackslash), \textbackslash(C(q,\textbackslash dot\{q\})\textbackslash), \textbackslash(G(q)\textbackslash) \\
\textbf{Constraint Handling} & Systematically incorporates state and control constraints & No built-in constraint handling \\
\textbf{Computation} & Solves a QP at every time step & Matrix-vector product; very fast per step \\
\textbf{Best For} & Systems with constraints, slow-medium dynamics & Fast manipulator control with known dynamics \\
\end{longtable}

\paragraph{Key Takeaway 1}\label{key-takeaway-1}}

MPC converts the control problem into a \textbf{repeated optimisation} problem.

\paragraph{Key Takeaway 2}\label{key-takeaway-2}}

The receding-horizon strategy provides \textbf{implicit feedback}.

\paragraph{Key Takeaway 3}\label{key-takeaway-3}}

MPC\textquotesingle s greatest strength is \textbf{constraint handling}.

\paragraph{Key Takeaway 4}\label{key-takeaway-4}}

Computed torque remains superior at \textbf{very high rates} (\textgreater{} 1 kHz).

\subsection{Next Week: Adaptive Control}\label{next-week-adaptive-control}}

\textbf{Building on Weeks 1 and 2, we will cover:}

\begin{itemize}
\tightlist
\item
  Parameter estimation and adaptive control for manipulators with unknown dynamics
\item
  Model Reference Adaptive Control (MRAC)
\item
  Linearity-in-parameters and the adaptive computed torque controller
\item
  Stability proofs for adaptive systems using Lyapunov-based methods
\end{itemize}

\subsubsection{Preparation}\label{preparation}}

\begin{itemize}
\tightlist
\item
  Complete all MATLAB activities from today and save your scripts
\item
  Review the linearity-in-parameters property from Week 1
\item
  Think about: \emph{"If the controller does not know \textbackslash(M(q)\textbackslash), \textbackslash(C(q,\textbackslash dot\{q\})\textbackslash), or \textbackslash(G(q)\textbackslash), how could it learn these online while still maintaining stability?"}
\end{itemize}

\textbf{Control for Robots}

University of West London \textbar{} Academic Year 2025--26

Week 2: Model Predictive Control (MPC)

This lecture material is for educational use only.
